%% sections/anchor.tex
\section{Three Belief-Update Plans}
\label{sec:anchor}

\subsection{Overview}

This section specifies the three \emph{plans} by which a utility
agent translates its prior and the messages it has received into
a posterior valuation each period. Every agent starts with the same
Lopez-Lira prior of equation~\eqref{eq:util_prior}; every agent
scores candidate orders with the same expected-utility rule of
equation~\eqref{eq:eu}; every plan feeds its output---a single
number $V_i^{\text{post}}$, in experimental cents---into that
score. The only thing that differs across plans is the function
that maps the prior plus peer messages into the posterior.

Plan~I is deterministic and runs closed-form. Plan~II and Plan~III
fire a parallel LLM call per utility agent at every period boundary,
and differ only in the prompt: Plan~II hands the model the
closed-form risk-utility expression for that agent's type, while
Plan~III names only the risk-preference category. A failed call in
Plan~II or Plan~III degrades silently to Plan~I for that agent in
that period, so every run is replayable even when the key is
absent and the mulberry32 PRNG alone drives the downstream dynamics.

\subsection{Plan I: algorithmic blend with experience-weighted trust}
\label{sec:plan1}

Let $V_i^{\text{prior}}$ denote the agent's Lopez-Lira prior from
equation~\eqref{eq:util_prior}, let $\mathcal{M}_i$ be the set of
claimed valuations received from foreign peers during the current
period (messages from the agent itself are excluded), and let
$k_i=\texttt{roundsPlayed}_i\in\{0,1,2,\ldots\}$ be the agent's
endogenous experience counter. Plan~I computes the period-end
posterior as the convex combination
\begin{equation}
V_i^{\text{post}} \;=\; w(k_i)\cdot V_i^{\text{prior}}
\;+\;\bigl(1-w(k_i)\bigr)\cdot\frac{1}{|\mathcal{M}_i|}\sum_{m\in\mathcal{M}_i}\hat v_m,
\qquad
w(k_i) \;=\; 0.6 + 0.1\,\min(3,\,k_i),
\label{eq:plan1_body}
\end{equation}
where the weight $w$ ramps deterministically with experience: a
novice ($k_i=0$) moves with a substantial peer contribution
($w=0.60$), a three-round veteran ($k_i\ge 3$) barely moves
($w=0.90$). The cap at $k_i=3$ matches the four-round session
structure: after three rounds of experience the agent is at its
terminal skepticism for the run. When $\mathcal{M}_i$ is empty---no
peer sent a valid claim during the period---Plan~I keeps the prior
unchanged, $V_i^{\text{post}}=V_i^{\text{prior}}$.

Plan~I is the offline baseline to which Plans~II and~III are compared.
Because it consults no external resource, it executes in constant
time per period and its output is a pure function of the seeded
PRNG plus the received-message set. Two runs with the same
(population, seed) pair therefore produce identical trajectories,
and the downstream \texttt{Market} replay can be reconstructed
frame-by-frame from the append-only history arrays without
re-running any agent.

\subsection{Plan II: LLM update with explicit utility forms}
\label{sec:plan2}

Plan~II replaces equation~\eqref{eq:plan1_body} with a live language
model call. Let $A_U$ be the set of utility agents in the current
population. At every period boundary, once the messaging round has
closed and the trust EMA has updated (Section~\ref{sec:agents}), the
engine snapshots each agent's received-message set onto the agent
and fires a single \texttt{chat/completions} call per agent in
parallel. The system prompt is fixed and instructs the model to
return a single number with no prose. The user prompt
$\pi_i^{\text{II}}$ encodes six facts:
\begin{enumerate}
\item the paper constants $(T,\E[d])$ and the current period index~$t$;
\item the agent's risk-preference category~$U_i\in\{\Uloving,\Uneutral,\Uaverse\}$;
\item the \emph{closed-form} expression of the risk-utility functional~$U_i$;
\item the agent's bias mode $b_i\in\{\text{over},\text{none},\text{under}\}$;
\item the agent's current mark-to-market wealth
$w_i = \text{cash}_i + \text{inventory}_i\cdot p^\star$ with
$p^\star=\max(\text{last}, \FV_t)$;
\item the set of peer claims $\{\hat v_m\}_{m\in\mathcal{M}_i}$ received during
the current period.
\end{enumerate}
The third item is the load-bearing element. Concretely the prompt
reads
\begin{quote}\small
You, agent $i$, are a \emph{risk} trader. Your expected utility over
a next-period wealth $w$ takes the closed form
\[
U_i(w) \;=\; (w/w_0)^2 \quad\text{/}\quad w/w_0 \quad\text{/}\quad \sqrt{w/w_0}
\]
for risk-loving, risk-neutral, or risk-averse, respectively. Your
current mark-to-market wealth is $w_i$. During the current period
your peers broadcast the following private valuations:
$\{\hat v_m\}$. Market: $T$ periods, expected dividend $\E[d]$,
current fundamental value $\FV_t$ cents. What is your private
subjective valuation per share for next period, in cents?
\end{quote}
where \emph{risk} is replaced by the agent's own category. The
reply is parsed for the first signed decimal, converted to a float,
clamped to a sanity interval, and written to the agent's
\texttt{llmBeliefs} slot in the engine context. On the next period's
first call to \texttt{updateBelief}, the slot is consumed in place
of equation~\eqref{eq:plan1_body}. If the clamp rejects the reply,
or the parse fails, or the network errors, the slot is left empty
and the agent falls through to Plan~I for that period. The next
period's LLM call then runs fresh, so a missed update never cascades.

The key feature of Plan~II is that the prompt explicitly contains
the functional form $(w/w_0)^2$, $w/w_0$, or $\sqrt{w/w_0}$ for
the agent's assigned risk preference. A language model that
already knows what a risk-averse agent looks like does not need
the formula; a language model that does not know (or that
translates ``risk-averse'' into a qualitatively different
curvature) behaves differently when the formula is present. The
contrast between Plan~II and Plan~III measures exactly that.

\subsection{Plan III: LLM update with risk label only}
\label{sec:plan3}

Plan~III uses the same period-boundary wiring as Plan~II with
identical timing, identical clamping, identical fallback, and the
same parallel batch of chat completions. The \emph{only} difference
is the user prompt: the closed-form expressions are deleted. The
prompt states
\begin{quote}\small
You, agent $i$, are a \emph{risk-loving} / \emph{risk-neutral} /
\emph{risk-averse} trader. Your current mark-to-market wealth is
$w_i$. During the current period your peers broadcast the following
private valuations: $\{\hat v_m\}$. Market: $T$ periods, expected
dividend $\E[d]$, current fundamental value $\FV_t$ cents. What is
your private subjective valuation per share for next period, in
cents?
\end{quote}
No $U_L / U_N / U_A$ formula is supplied and no reference is made
to expected utility; the model must translate the risk label into
an implicit curvature internally. Everything else---the mark-to-market
wealth, the peer claims, the market constants, the prompt preamble,
and the single-number return requirement---is byte-for-byte
identical to Plan~II.

If Plan~II and Plan~III produce measurably different distributions
over $V_i^{\text{post}}$ or over downstream market-quality
statistics, the difference is attributable to prompt structure
alone. If they agree, the LLM in use has a stable internal
representation of the three risk archetypes that is robust to
whether the functional form is named in the prompt.

\subsection{Consume-and-cache on the period boundary}

The LLM call in Plan~II and Plan~III is fired \emph{after} the
period's messaging round closes but \emph{before} the next period's
first decision tick, so that the prompt can include the peer
claims observed during the period that just ended. Concretely the
engine:
\begin{enumerate}
\item copies each agent's \texttt{receivedMsgs} array into a
snapshot field, so the downstream call reads a stable message set;
\item fires one \texttt{chat/completions} call per agent in a single
\texttt{Promise.all} batch, so network latency is one round trip
rather than $|A_U|$ round trips;
\item merges each successful reply into a context-level
\texttt{llmBeliefs} map keyed by agent id;
\item on the next period's first call to \texttt{updateBelief},
consumes the entry in place of equation~\eqref{eq:plan1_body} and
clears the map slot so a second call in the same period reverts to
Plan~I.
\end{enumerate}
This \emph{consume-and-cache} pattern keeps the tick loop synchronous
and non-blocking: agents continue to quote and trade at the usual
tick cadence while the LLM batch runs in the background, and the
result is applied on the very next period boundary. If the call is
still in flight when the next boundary arrives, the
\texttt{llmBeliefs} slot is empty and the agent runs Plan~I for that
period, so the engine never stalls waiting for a network response.

\subsection{Determinism and reproducibility}

Plan~I is fully deterministic under the seeded mulberry32 PRNG: two
runs with the same (population, seed) pair produce byte-identical
trace, snapshot, and event logs. Plan~II and Plan~III introduce one
source of non-determinism---the LLM reply---and the engine handles
it by keeping the fallback path closed-form. If the same seed is
run twice under Plan~II, the two runs will in general disagree
wherever the language model's reply differs between calls, but
both runs will produce valid simulations and both will respect the
clamp~and the parse. Reporting therefore always pairs a plan label
with a seed when comparing statistics across runs.

A missing reply, a parse failure, or an HTTP error silently drops
the update for that agent in that period, and the agent falls
through to Plan~I. This preserves the invariant that every run can
be reproduced without network access---Plan~II with zero successful
calls is definitionally identical to Plan~I---so the simulator can
be used as a pedagogical or offline replay tool even when the API
endpoint is unreachable.
